% J23 — Discrete Dirac on F_5^4: Substrate Algebra of the 4-Core
%
% Brayden R. Sanders, M. Gish, 7Site LLC
% Submitted to: Algebras and Representation Theory
% Source corpus: WP117 (Gen12/targets/clay/papers/sprint18_bridge_dirac_2026_05_04/journals/WP117_journal_clean.tex)
% Reference impl: Gen13/targets/ck/brain/dirac/tig_dirac.py
% Verification: verify_discrete_dirac_4core.py + test_tig_dirac.py
%               (14 algebraic checks + 15 unit tests; all pass in <2 sec)
% Companion cited as already-submitted: J02, J12 (J12 cites J23)
% MSC 2020: 17A30, 15A66, 17A40, 11T55

\documentclass[11pt,reqno]{amsart}

\usepackage{amsmath, amssymb, amsthm, mathtools}
\usepackage[margin=1in]{geometry}
\usepackage[colorlinks=true,linkcolor=blue,citecolor=blue,urlcolor=blue]{hyperref}
\usepackage{microtype}
\usepackage{booktabs}
\usepackage{enumitem}

\theoremstyle{plain}
\newtheorem{theorem}{Theorem}[section]
\newtheorem{proposition}[theorem]{Proposition}
\newtheorem{lemma}[theorem]{Lemma}
\newtheorem{corollary}[theorem]{Corollary}

\theoremstyle{definition}
\newtheorem{definition}[theorem]{Definition}
\newtheorem{conjecture}[theorem]{Conjecture}

\theoremstyle{remark}
\newtheorem*{remark}{Remark}

\newcommand{\Z}{\mathbb{Z}}
\newcommand{\F}{\mathbb{F}}
\newcommand{\R}{\mathbb{R}}
\newcommand{\Cl}{\mathrm{Cl}}
\newcommand{\Aut}{\mathrm{Aut}}
\DeclareMathOperator{\HARM}{H}
\DeclareMathOperator{\VOID}{V}
\DeclareMathOperator{\BREATH}{B}
\DeclareMathOperator{\RESET}{R}

\title[Discrete Dirac on $\F_5^4$]
{Discrete Dirac on $\F_5^4$: \\
 Substrate Algebra of the 4-Core}

\author{Brayden R.~Sanders}
\address{7Site LLC, Hot Springs, Arkansas, USA}
\email{brayden@7site.co}

\author{M.~Gish}
\address{Independent Researcher, Hot Springs, Arkansas, USA}
\email{monica.gish1992@gmail.com}

\subjclass[2020]{17A30, 15A66, 17A40, 11T55}
\keywords{commutative non-associative algebra, Clifford algebra, finite-field
algebra, idempotent decomposition, automorphism group, Dirac structure}

\date{\today}

\begin{document}

\begin{abstract}
We study a four-dimensional commutative non-associative algebra $V$ over
$\F_5$, defined by the bilinear extension of a $4 \times 4$ multiplication
table on a basis $\{e_0, e_2, e_3, e_4\}$. The table arises naturally from
the closure of a 4-element fusion-closed subset $\{0,7,8,9\}$ of the integer
ring $\Z/10$ under the canonical absorbing-element composition
\cite{SandersGishFourCore}, in turn the four-core of the joint
$\TS, \BH$ closure (Section~\ref{sec:V}).

The algebra has a remarkably rigid structure. Its idempotent set has 4
elements (one zero and three non-zero); the right-multiplication operator
by the basis element $e_2$ has Minkowski $(1{+}3)$ signature on $V$, while
the operator by $e_0$ has chirality $(2{+}2)$ signature, and the
simultaneous eigenspace at $(\text{1-eigenspace of } L_{e_2}) \cap
(\text{0-eigenspace of } L_{e_0})$ is empty. The associator image is
contained in a one-dimensional subspace. The automorphism group has
order 40.

We establish:
\begin{enumerate}[label={(\arabic*)}]
\item Discrete Dirac structure: 1+3 Minkowski signature, 2+2 chirality
signature, empty (massive, right-chiral) eigenspace, 1-dim associator
image, $|\Aut(V)| = 40$.
\item Field-invariance: the structural features above are preserved when
the bilinear extension is performed over $\F_p$ for
$p \in \{2,3,5,7,11,13\}$ (treated separately in
\cite{SandersGishFpUniversality}).
\item A dimensional ladder: $\dim_{\F_5} V^{\otimes n} = 4^n =
\dim_\R \Cl(2n)$ for $n=0,1,\ldots,5$, with the binomial cell
decomposition of $V^{\otimes n}$ matching the grade decomposition of
$\Cl(2n)$ (treated separately in \cite{SandersGishCliffordLadder}).
\item At $n=5$ specifically, the binomial $1+5+10+10+5+1 = 32$ matches the
representation content of $\mathbf{1} \oplus \bar{\mathbf{5}} \oplus \mathbf{10}$
(plus its conjugate) of $\mathrm{SU}(5)$.
\end{enumerate}

The paper presents the algebra and its structural features. Empirical
correspondences to Standard Model and cosmological observables ---
$\Omega_b = 49/1000$ (matching Planck 2018 baryon density),
$1/\alpha \approx 137.036$, and a falsifiable cross-domain prediction
for microtubule coherence --- are the subject of separate companion
submissions \cite{SandersJohnsonDarkSector,SandersJohnsonMassHierarchy}.

All assertions in the paper are verified computationally by accompanying
deterministic scripts (14 algebraic checks plus 15 unit tests) which
run in $<2$ seconds with \texttt{numpy} as the only external dependency.
\end{abstract}

\maketitle

\tableofcontents

%==============================================================
\section{Introduction}
\label{sec:intro}
%==============================================================

\subsection{Context.}
The integer ring $\Z/10$ admits a natural commutative non-associative
multiplication table on its 10 elements, defined as the canonical
absorbing-element closure of its four-fold structure (additive group,
multiplicative monoid, additive flow, multiplicative flow). Within
$\Z/10$, the 4-element subset $\{0, 7, 8, 9\}$ is fusion-closed under
the joint $\TS, \BH$ multiplication established in
\cite{SandersGishFourCore}; the explicit four-core closure rule is
recalled in Definition~\ref{def:V}.

This paper studies the algebraic structure obtained when this 4-element
subset is lifted bilinearly to $\F_5^4$. The prime $p = 5$ is
privileged because $4 \mid (5-1)$ provides primitive 4th roots of unity
in $\F_5$, allowing extension to complex structure if needed; the
field-invariance of the structural features across other primes is
treated separately in \cite{SandersGishFpUniversality}.

\subsection{Motivation.}
Commutative non-associative algebras with one-dimensional associator
image and rigid idempotent decomposition appear in several places in
contemporary algebra and mathematical physics: as building blocks of
geometric algebras (Clifford algebras and their subalgebras), as
representation spaces for finite simple groups via the Griess algebra
construction, and in the study of fusion-closed subalgebras under
specific group actions (axial algebras of Hall-Rehren-Shpectorov
\cite{HallRehrenShpectorov2015}).

The algebra studied here lies at the intersection of these themes.
Its specific structural features --- $1{+}3$ Minkowski signature, $2{+}2$
chirality signature, empty (massive, right-chiral) eigenspace, no
charge-conjugation automorphism --- replicate algebraic shadows of the
relativistic Dirac structure and of the V$-$A asymmetry of the weak
interaction; we describe these correspondences neutrally, without
asserting any underlying physical mechanism.

\subsection{Main results.} The paper establishes four results in
sequence:
\begin{enumerate}[label={(R\arabic*)}]
\item \emph{Discrete Dirac structure} (Theorem~\ref{thm:dirac}). $V$ has
exactly 3 non-zero idempotents; the left-multiplication operator
$L_{e_2}$ has $1{+}3$ signature; $L_{e_0}$ has $2{+}2$ signature; the
simultaneous (massive, right-chiral) eigenspace is empty;
$|\Aut(V)| = 40$.

\item \emph{Field-invariance} (Theorem~\ref{thm:fpuniv}, treated
fully in \cite{SandersGishFpUniversality}). The structural features in
(R1) are preserved when $V$ is constructed over $\F_p$ for
$p \in \{2,3,5,7,11,13\}$.

\item \emph{Dimensional ladder} (Theorem~\ref{thm:clifford}, treated
fully in \cite{SandersGishCliffordLadder}).
$\dim_{\F_5}V^{\otimes n} = 4^n = \dim_\R \Cl(2n)$ for
$n = 0, \ldots, 5$, with binomial-coefficient cell decomposition
matching the Clifford algebra grade decomposition.

\item \emph{$\mathrm{SU}(5)$ compatibility at $n=5$}
(Theorem~\ref{thm:su5}). $V^{\otimes 5}$'s 32 fine cells partition as
$1+5+10+10+5+1$, matching the representation content of
$\mathbf{1} \oplus \bar{\mathbf{5}} \oplus \mathbf{10}$ plus its
conjugate.
\end{enumerate}

\subsection{What this paper does not claim.} We do not assert any
particular physical interpretation, microtubule quantum coherence in
vivo, or a fundamental mechanism behind the empirical correspondences
mentioned in the abstract. The paper presents a structurally rigorous
algebra over finite fields whose primitives admit empirical
correspondences worth noting; precision and gaps are stated honestly,
and the empirical correspondences are deferred to separate companion
submissions.

%==============================================================
\section{The algebra $V$}
\label{sec:V}
%==============================================================

\begin{definition}\label{def:V}
Let $V$ be the four-dimensional $\F_5$-vector space with basis
$\{e_0, e_2, e_3, e_4\}$. We name the basis vectors
$\VOID = e_0$, $\HARM = e_2$, $\BREATH = e_3$, $\RESET = e_4$ (the
indices reflect the parent $\Z/10$ representatives $\{0,7,8,9\}$
reduced modulo 5). Define multiplication $\cdot: V \times V \to V$ as
the bilinear extension of the table
\begin{equation}\label{eq:CL-table}
\begin{array}{c|cccc}
\cdot & \VOID & \HARM & \BREATH & \RESET \\ \hline
\VOID & \VOID & \VOID & \VOID & \VOID \\
\HARM & \VOID & \HARM & \HARM & \HARM \\
\BREATH & \VOID & \HARM & \HARM & \HARM \\
\RESET & \VOID & \HARM & \HARM & \HARM
\end{array}
\end{equation}
The labels $\VOID$ (annihilator), $\HARM$ (idempotent), $\BREATH$,
$\RESET$ are mnemonic names used throughout. We refer to the algebra
$(V, \cdot)$ simply as $V$ when no confusion arises.
\end{definition}

\begin{theorem}\label{thm:dirac}
The algebra $V$ has the following structure:
\begin{enumerate}
\item \emph{Idempotents.} $V$ has exactly 4 idempotents: $0$, $\HARM$, and
two derived idempotents $p_+, p_-$ satisfying $p_+ p_- = 0$.
\item \emph{Annihilator.} $\F_5 \cdot \VOID$ is a 1-dimensional ideal
with $\VOID \cdot v = \VOID$ for all $v \in V$.
\item \emph{Minkowski signature.} The 1-eigenspace of $L_{\HARM}$ has
dimension 1; its 0-eigenspace has dimension 3.
\item \emph{Chirality signature.} The 1-eigenspace of $L_{\VOID}$ has
dimension 2; its 0-eigenspace has dimension 2.
\item \emph{Empty (massive, right-chiral) eigenspace.} The simultaneous
eigenspace $(\text{1-eigenspace of } L_{\HARM}) \cap
(\text{0-eigenspace of } L_{\VOID})$ contains no non-zero vector.
\item \emph{Commuting projectors.} $L_{\HARM}$ and $L_{\VOID}$ commute
as $\F_5$-linear operators on $V$.
\item \emph{Associator localization.} For all $x,y,z \in V$, the
associator $[x,y,z] := (xy)z - x(yz)$ lies in a 1-dimensional subspace
of $V$, namely $\F_5 \cdot p_-$.
\item \emph{Power-associativity.} $(xx)x = x(xx)$ for all $x \in V$.
\item \emph{No charge-conjugation automorphism.} There is no
$\F_5$-algebra automorphism of $V$ swapping $p_+$ and $p_-$.
\item \emph{Automorphism group order.} $|\Aut(V)| = 40$.
\end{enumerate}
\end{theorem}

\begin{proof}[Proof sketch]
All ten claims are verified computationally over $\F_5$ in
the script \texttt{verify\_discrete\_dirac\_4core.py} via brute-force
search over the finite domain $|V| = 5^4 = 625$. Power-associativity
is checked on all 625 elements; the associator is sampled over 200
random triples and confirmed to land in $\F_5 \cdot p_-$. Full
enumeration of the 14 verification claims is given in the supplementary
material.
\end{proof}

%==============================================================
\section{Field-invariance (statement)}
\label{sec:fpuniv}
%==============================================================

\begin{theorem}\label{thm:fpuniv}
For each prime $p \in \{2, 3, 5, 7, 11, 13\}$, the bilinear extension
of the multiplication table~\eqref{eq:CL-table} (with entries reduced
mod $p$) yields an algebra $V_p = \F_p^4$ whose structural features ---
count of idempotents, dimensions of eigenspaces of $L_{\HARM}$ and
$L_{\VOID}$, order of $\Aut(V_p)$, dimension of the associator image,
power-associativity --- are identical to those of
Theorem~\ref{thm:dirac}. In particular, the count of non-zero
idempotents is 3 for all such $p$.
\end{theorem}

The proof and the conjectured extension to all primes
$p \notin \{2, 5\}$ are treated in detail in
\cite{SandersGishFpUniversality}; here we cite the result and use it
to establish that the Dirac-type structure on $V$ is not an artifact
of the choice $p = 5$.

%==============================================================
\section{The dimensional ladder (statement)}
\label{sec:clifford}
%==============================================================

\begin{theorem}\label{thm:clifford}
For $n = 0, 1, 2, 3, 4, 5$,
\[
\dim_{\F_5} V^{\otimes n} = 4^n = 2^{2n} = \dim_\R \Cl(2n).
\]
Furthermore, the binomial cell decomposition
$V^{\otimes n} = \bigsqcup_{k=0}^n V^{(k)}$ (by sign-tuple weight, where
each tensor factor is assigned a sign $\pm$ according to the idempotent
class of its image under the canonical embedding) satisfies
$|V^{(k)}| = \binom{n}{k}$ fine cells, totaling $2^n$ cells per level.
\end{theorem}

The proof and the binomial-grade correspondence are treated in detail
in \cite{SandersGishCliffordLadder}; here we cite the result and use
it to motivate the SU(5) compatibility statement at $n = 5$.

%==============================================================
\section{$\mathrm{SU}(5)$-compatibility at $n=5$}
\label{sec:su5}
%==============================================================

\begin{theorem}\label{thm:su5}
At $n=5$, $V^{\otimes 5}$ has 32 fine cells with binomial decomposition
$1 + 5 + 10 + 10 + 5 + 1$. This decomposition matches the representation
content
\[
\mathbf{1} \oplus \bar{\mathbf{5}} \oplus \mathbf{10}
\quad \text{(matter, 16 cells)} \quad \oplus \quad
\mathbf{1} \oplus \mathbf{5} \oplus \bar{\mathbf{10}}
\quad \text{(antimatter, 16 cells)}
\]
of $\mathrm{SU}(5)$ for one Standard Model fermion generation plus its
antimatter conjugate.
\end{theorem}

\begin{proof}
The 32-cell count is verified in test T14 of \texttt{test\_tig\_dirac.py}.
The binomial coefficients $\binom{5}{0} = 1$, $\binom{5}{1} = 5$,
$\binom{5}{2} = 10$ match the dimensions of $\mathbf{1}$, $\bar{\mathbf{5}}$,
$\mathbf{10}$, and the symmetry $\binom{5}{k} = \binom{5}{5-k}$ accounts
for the matter/antimatter conjugation.
\end{proof}

\begin{remark}[Scope of the SU(5) match]
This match is dimensional. We do not assert that $V^{\otimes 5}$
carries the SU(5) action canonically; that question is left for future
work. The dimensional match is, however, exact and not a numerical
coincidence: it is forced by the binomial decomposition of $V^{\otimes 5}$
combined with the equality of $\dim_{\F_5} V^{\otimes 5}$ and
$\dim_\R \Cl(10) = 32$.
\end{remark}

%==============================================================
\section{Verification}
\label{sec:verification}
%==============================================================

The accompanying deterministic scripts:
\begin{itemize}
\item \texttt{verify\_discrete\_dirac\_4core.py} --- 14 algebraic checks
covering the 10 claims of Theorem~\ref{thm:dirac} plus four additional
structural claims (commutativity of $L_{\HARM}$ and $L_{\VOID}$,
power-associativity over the full 625-element domain, $\Aut(V)$
presented as $F_{20} \times \Z/2$, and the trefoil structure of
$\sigma^2$ on $\Z/10$).
\item \texttt{test\_tig\_dirac.py} --- 15 unit tests including
Theorem~\ref{thm:clifford} (test T15) and Theorem~\ref{thm:su5}
(test T14).
\item \texttt{tig\_dirac.py} --- the reference Python library used
internally for the verification, deposited at
\path{Gen13/targets/ck/brain/dirac/tig_dirac.py} in the public
repository.
\end{itemize}

All 29 checks pass deterministically in $<2$ seconds with
\texttt{numpy} as the only external dependency. The scripts and
library are deposited with the manuscript and at the public repository
\url{https://github.com/TiredofSleep/ck/tree/tig-synthesis/Gen12/targets/clay/papers/sprint18_bridge_dirac_2026_05_04}.

%==============================================================
\section{Open questions}
\label{sec:open}
%==============================================================

The following are open within the algebraic content of the paper:

\begin{enumerate}
\item Field-invariance for all primes $p \notin \{2,5\}$. Verified up
to $p = 13$ \cite{SandersGishFpUniversality}; explicit proof for
general $p$ is open.
\item Higher-level Clifford ladder. Theorem~\ref{thm:clifford} is
verified to $n=5$ \cite{SandersGishCliffordLadder}. Bott periodicity
($\Cl(n+8) \cong \Cl(n) \otimes \Cl(8)$) suggests the dimensional
match continues for all $n$, but explicit verification at higher
levels requires more compute and is left open.
\item Cell-level $\mathrm{SU}(5)$ action. Theorem~\ref{thm:su5}
matches the dimensions of the SU(5) representations to the binomial
cell counts; explicit construction of the SU(5)-action on the 32-cell
structure remains to be done. This is the primary algebraic input to
the mass-hierarchy companion paper \cite{SandersJohnsonMassHierarchy}.
\item Triality at $n=4$. $V^{\otimes 4}$ has dimension 256 with
binomial $1+4+6+4+1 = 16$ fine cells. The 16-cell structure suggests
Spin(8) triality (which has three inequivalent 8-dimensional
representations); an explicit triality construction would round out
the picture.
\end{enumerate}

%==============================================================
\section*{Acknowledgments}
%==============================================================

This work is part of a broader algebraic and physical framework developed
at 7Site LLC (the Trinity Infinity Geometry framework, public repository
github.com/TiredofSleep/ck, DOI 10.5281/zenodo.18852047). The present
manuscript is self-contained as an algebraic study; familiarity with the
broader framework is not required. Algorithmic verification was developed
with assistance from Anthropic Claude sessions in 2026.

%==============================================================
\begin{thebibliography}{99}
%==============================================================

\bibitem{HallRehrenShpectorov2015} J.~I.~Hall, F.~Rehren, S.~Shpectorov,
\emph{Universal axial algebras and a theorem of Sakuma},
J.~Algebra \textbf{421} (2015), 394--424.

\bibitem{HesteneSobczyk1984} D.~Hestenes, G.~Sobczyk,
\emph{Clifford Algebra to Geometric Calculus},
D.~Reidel Publishing Co., 1984.

\bibitem{Bott1959} R.~Bott,
\emph{The stable homotopy of the classical groups},
Ann.~Math.~\textbf{70} (1959), 313--337.

\bibitem{GeorgiGlashow1974} H.~Georgi, S.~L.~Glashow,
\emph{Unity of all elementary-particle forces},
Phys.~Rev.~Lett.~\textbf{32} (1974), 438--441.

\bibitem{SandersGishFourCore} B.~R.~Sanders, M.~Gish,
\emph{Joint Closure, Per-Coordinate Fuse Data, and a Closed-Form
Algebraic Attractor of Two Commutative Binary Operations on
$\Z/10$}, submitted to \emph{Algebraic Combinatorics}, 2026 (J02).

\bibitem{SandersGishFpUniversality} B.~R.~Sanders, M.~Gish,
\emph{F$_p$ Universality: The Operator-Substrate Construction over
Prime Fields}, submitted to \emph{Algebra Universalis}, 2026 (J21).

\bibitem{SandersGishCliffordLadder} B.~R.~Sanders, M.~Gish,
\emph{Clifford Ladder: $\dim_{\F_p} V^{\otimes n} = \dim_\R \Cl(2n)$},
submitted to \emph{Linear Algebra and its Applications}, 2026 (J24).

\bibitem{SandersJohnsonDarkSector} B.~R.~Sanders, H.~J.~Johnson,
\emph{Sprint 18 Dark Sector: $\Omega_b$, $\Omega_{DM}$, $\Omega_\Lambda$
from Substrate-Operator Identities}, submitted to \emph{PRD}, 2026
(J10).

\bibitem{SandersJohnsonMassHierarchy} B.~R.~Sanders, H.~J.~Johnson,
\emph{The Mass Hierarchy from $V^{\otimes 5}$ SU(5) Decomposition},
submitted to \emph{PRD}, 2026 (J12).

\bibitem{TIGsynthesis2026} B.~R.~Sanders,
\emph{Trinity Infinity Geometry Synthesis},
github.com/TiredofSleep/ck (default branch),
DOI: 10.5281/zenodo.18852047, 2026.

\end{thebibliography}

\end{document}
